\section[曲线上主 $G$-丛的模叠]{曲线上主 $G$-丛的模叠（3 月 30 日）}

\srcpage{46}

我们希望把前几讲见过的阿贝尔对偶推广到 $G$ 非阿贝尔的情形。一个早期的困难是：$\C[G]$ 非交换，因此没有显然的对偶 $\Spec(\C[G],*)$。为了绕开这个问题、仍得到对偶，我们需要另找一个交换的对象。我们将用拓扑场论来做到这一点。

这段故事可以从两端讲起：从问题出发，或从解答（TFT）出发；我们从问题开始。具体地说，我们要把类域论（class field theory）中出现的 $\Pic(C)$ 换成一个非阿贝尔的对应物：自守形式（automorphic form）的理论。

设 $C$ 是域 $k$ 上的光滑紧合曲线，$k$ 或为 $\F_q$，或为 $\C$（后一情形 $C$ 是 Riemann 面）。在阿贝尔情形，我们研究的是 $\Pic(C)$ 上的函数（$k=\F_q$）或层（更一般的 $k$，特别地 $k=\C$）；$\Pic(C)$ 是 $C$ 上线丛的模空间。现在把这里的「线丛」换成向量丛，或更一般地，换成主 $G$-丛。

设 $G$ 是 $k$ 上的代数群，仿射，并且（通常）约化。我们暂不定义「约化」，只说它包括 $\mathrm{GL}_n$、$\mathrm{SO}_n$、$\mathrm{Sp}_n$ 这些常客。主 $G$-丛（principal $G$-bundle）是一个空间 $P$，带有映射 $P\to C$，以及 $C$ 上的右 $G$-作用；并且局部于 $C$ 有 $P|_U\cong U\times G$，其中 $G$ 作用于 $G$ 而不作用于 $U$。当 $k=\F_q$ 时，必须说清「局部」的含义：我们指 étale 局部。对 $G=\mathrm{GL}_n$，这与秩 $n$ 向量丛是等价的数据。空间上的主丛构成群胚：若仿照向量丛之间映射的定义写出主丛之间映射的概念，你会发现所有这样的映射都是同构。特别地，主丛有自同构。

令 $\Bun_G(C)$ 为 $C$ 上主 $G$-丛的模叠（moduli stack）。$\Bun_G(C)$ 的函子把环 $R$ 送到 $C\times\Spec R$ 上主 $G$-丛的群胚。当 $G=\mathrm{GL}_1=\G_m$ 时，有等价

\srcpage{47}

\[\Bun_{\G_m}(C)\cong \ul{\Pic}(C):=\Pic(C)\times pt/\G_m.\]

也就是说，主 $\G_m$-丛在同构意义下与线丛是等价的数据，但所有主 $\G_m$-丛的自同构群都是 $\G_m$。

设 $k=\F_q$；扔掉 $\Bun_G(C)(\F_q)$ 中的稳定化子，便得到主 $G$-丛同构类的离散集。它一般并不有限，但没有任何几何。把这个离散集记作 $|\Bun_G(C)(\F_q)|$。\yizhu{「扔掉稳定化子」即取群胚的底层同构类集合；原文以此强调该集合不带几何结构。}{A（原文 p.47）}

\begin{definition}
$C$ 上 $G$ 的（无分歧）自守函数（automorphic function）是 $|\Bun_G(C)(\F_q)|$ 上的 $\C$ 值函数。\footnote{若愿意，也可以取 $\Ql$ 值函数。}
\end{definition}

一般地，当这个集合无限时，我们会对这些函数加某种增长条件，例如紧支。

阿贝尔情形中，我们给出了 $\Pic(C)$ 上函数的 Fourier 分解；现在我们要对无分歧自守形式做谱分解。我们可以使用 Weil 对 $|\Bun_G(C)(\F_q)|$ 的描述，它（相对地！）具体，并且类似于我们对 $\Pic(C)$ 的描述。

首先，$C$ 上任何主 $G$-丛在一般点处都平凡。要看出这一点需要想一下；展开来说，就是要找 $P\to C$ 的有理截面，即在 $C$ 去掉有限多个点上的截面。而带有截面的主丛可以平凡化。

其次，任何 $G$-丛在一点的形式邻域上平凡。于是我们可以在几乎处处平凡化，即在 $C$ 去掉有限多个点 $x_i$ 后平凡化；再在每个 $x_i$ 的形式邻域中平凡化；然后描述这些描述之间的转移函数。这里需要一个重要工具——Beauville–Laszlo 定理，它说我们确实可以由形式的转移数据粘出一个 $G$-丛。

对每个 $x_i$，有 Taylor 级数的完备局部环 $\oo_{x_i}$，它包含在 Laurent 级数的局部域 $K_{x_i}$ 中。在 $\Spec(\oo_{x_i})$ 与 $C\setminus x$ 的交叠上，得到两个平凡化的差，它落在 $G(K_{x_i})$ 中。这些数据汇成限制乘积（restricted product）中的一个元素

\[\prod_{x\in C}{}^{\prime} G(K_x),\]

意思是除有限多个外全为零。Weil 随后确定了这些陪集。\yizhu{原文用 $F_C$ 记 $C$ 的函数域，用 $\A_C$ 记其 adele 环；下文照此。}{A（原文 p.47）}

\begin{theorem}[Weil]
设 $k=\F_q$。则
\begin{align*}
|\Bun_G(C)(k)| &\cong G(F_C)\backslash \prod_{x\in C}{}^{\prime} G(K_x)\Big/ \prod_{x\in C}{}^{\prime} G(\oo_x)\\
&\cong G(F_C)\backslash G(\A_C)/G(\oo_{\A_C}).
\end{align*}
\end{theorem}

回忆阿贝尔情形也有一个形状类似的双陪集分解——主要差别在于：阿贝尔群对子群的商总是群，而这里未必。阿贝尔情形的局部因子

\[\mathrm{GL}_1(K_x)/\mathrm{GL}_1(\oo_x)=K_x^\times/\oo_x^\times\cong \Z\]

正是「除子即线丛」这一描述的来源。但在非阿贝尔情形，我们要对付更困难的东西：

\[\Gr_G:=G(K_x)/G(\oo_x),\]

称为 $G$ 的仿射 Grassmann 簇（affine Grassmannian）。它不是群：$G(\oo_x)$ 在 $G(K_x)$ 中不正规。你可以把 $\Gr_G$ 想成这样的数据空间：$C$ 上一个主 $G$-丛，连同 $C\setminus x$ 上的一个平凡化。

上述描述在数域情形有很好的类比。这里 $G$ 对应 idele 类群，$F/\Q$ 是有限扩张（大部分一般特征在 $F=\Q$ 时已经出现）。考虑对象

\[G(F)\backslash G(\A_F)/G(\oo_{\A_F}),\]

其中 $\A_F$ 是 adele 环，即 $\oo_F$ 在所有位点的完备化的乘积。对 $\Q$，$\A_\Q$ 是各 $\Q_p$ 与 $\R$ 的限制乘积。$\Q_p$ 这些因子称为「有限 adele」，它同构于 $\hat\Z\otimes_\Z\Q$。函数环是 $\prod_p \Z_p=\hat\Z$。也就是说，adele 有一个「熟悉的」有限部分（$\Q_p$）和一个「奇怪的」\srcpage{48}阿基米德部分（$\R$）——如果你不是数论出身，这可能和你习惯的相反。

于是 $G(\A_\Q)$ 是 $G(\A_{\mathrm{fin}})$ 与实 Lie 群 $G(\R)$ 的乘积。我们商掉 $G(\Q)$，它的好处是带有离散拓扑——这是 adele 的一个便利事实。于是有

\[G(\Q)\backslash G(\A_\Q)/\big(G(\hat\Z)\times K_\infty\big),\]

其中 $K_\infty$ 是 $G(\R)$ 的紧子群，例如其极大紧子群。上述空间正是这个类比所提示的 $\Bun_G(\Q)$，尽管它其实不那么拓扑。

空间

\[G(\Q)\cap \prod_p G(\Z_p)\,G(\R)/K_\infty\]

落在上述空间里。左边的交是分母中不含任何 $p$ 的幂的有理数，所以它就是 $G(\Z)$。也就是说，我们有 $G(\Z)\backslash G(\R)/K_\infty$；而断言是：这是一个连通分支，是有限多个连通分支之一；对某些特别好的群（例如单连通的约化群），这还是双射！

要点在于，这相对具体，尽管它偏袒一个位点（$\R$）而轻视其余位点。它把模形式理论与自守形式理论联系起来。

\begin{example}
设 $G=\mathrm{GL}_2$，此时这联系到古典得多的东西。我们大致在问：$\Spec(\Z)$ 上秩 $2$ 向量丛的模空间是什么，或者 $\Spec(\Z)$ 的某种完备化上又是什么？

非常天真地说，这应该像是秩 $2$ 自由 $\Z$-模，但在无穷远有一个平凡化。令 $\cB$ 为 $\C$ 的 $\R$-基构成的空间，它可以等同于同构 $\Z^2\otimes\R\to\R^2$ 的空间；固定一组称心的基便给出 $\cB\cong \mathrm{GL}_2(\R)$。

我们可以具体地把 $\cB$ 描述为 $\C^\times\times\C\setminus\R$：先选第一个基向量，再选第二个。若 $\cL$ 是 $\C$ 中格构成的空间，即带有同构 $\Lambda\otimes\R\cong\C$ 的自由秩 $2$ $\Z$-模 $\Lambda$，则 $\cL$ 接近基的空间 $\cB$，但我们把基在可逆整值变换下等同起来：

\[\cL\cong \mathrm{GL}_2(\Z)\backslash\cB.\]

因为我们用的是 $\C$，这个空间有自然的复结构。这一点不能推广到更高秩。

此事有若干变体，例如 $\cL(N)$：格，并且还固定一个模 $N$ 的基，即一个等同 $\cL/N\cong(\Z/N)^{\oplus 2}$。你可以把它想成秩 $2$ 丛连同 $(N)\subset\Spec\Z$ 上的平凡化。这里用到了重数：例如 $(p)$ 与 $(p^2)$ 不同，这是在点 $p\in\Spec\Z$ 的更高阶邻域上做平凡化。

另一个变体是 $\ol{\cL}$，即格模去 $\C$-位似（homothety）\footnote{练习：「homothety」怎么读？}，位似即重新缩放。这个空间是

\[\ol{\cL}\cong \mathrm{GL}_2(\Z)\backslash\cB/\C^\times\cong \mathrm{PSL}_2(\Z)\backslash\mathbb{H}.\]

下面这个引理告诉我们，可以在 adele 的语言中识别格的空间，但不带无穷远的因子。

\begin{lemma}
\[\cL\cong \mathrm{GL}_2(\Z)\backslash \mathrm{GL}_2(\R)\cong \mathrm{GL}_2(\Q)\backslash \mathrm{GL}_2(\A)/\mathrm{GL}_2(\oo_\A).\]
\end{lemma}

这再次用到 $\mathrm{GL}_2(\Z)=\mathrm{GL}_2(\Q)\cap \mathrm{GL}_2(\hat\Z)$；还需要 $\mathrm{GL}_2(\A_{\mathrm{fin}})$ 由 $\mathrm{GL}_2(\oo_\A)$ 与 $\mathrm{GL}_2(\Q)$ 生成。这让我们能从这种平等对待所有位点的自守描述，过渡到挑出 $\R$ 的实齐性空间描述。商掉 $\C^\times$ 也可以改为商掉极大紧子群 $\mathrm{O}_2$。

模形式正是在这里古典地出现：取 $\mathrm{PSL}_2(\Z)\backslash\mathbb{H}$ 上线丛的全纯截面（尽管还带增长条件）。它们包含在 $\mathrm{GL}_2(\Z)\backslash\cL$ 上的函数空间中。因此我们可以写下 $\Q$ 上 $\mathrm{GL}_2$ 的自守形式的定义：它是

\[\mathrm{GL}_2(\Q)\backslash \mathrm{GL}_2(\A)/\mathrm{GL}_2(\oo_\A)\]

上的函数。

还有 level $N$ 的版本，把 $\mathrm{GL}_2(\oo_\A)$ 换成 $\mathrm{GL}_2(\oo_\A^{(N)})$。但先别管这个，回到 level $1$ 的版本。我们要求函数在 $\mathrm{GL}_2(\R)$ 上光滑，并在无穷远位点（$\R$）有某种增长条件。
\end{example}

\srcpage{49}

对其他群也有类似做法。例如 $G=\mathrm{O}_n$ 时，构造正交丛的模空间：固定一个偶格 $\Lambda$，考察带有二次函数 $q:V\to\Z$ 的自由 $\Z$-模 $V$。给定这些数据，可以谈论 $V$ 的自同构何时正交，即何时保持 $q$，于是得到正交群 $\mathrm{O}(V,q)$。从而有双商

\[\mathrm{O}_q(\Q)\backslash \mathrm{O}_q(\A)/\mathrm{O}_q(\oo_\A).\]

这有非常古典的含义：$q$ 所属的属（genus）中的二次型同构类，意思是它们与 $q$ 在每个模 $N$ 下都等同。$\Theta$ 级数之类就是这些二次型集合上的某些函数。

对任意数域 $F$ 上的任意群 $G$，我们都能玩这个游戏，把「$\Bun_G$」定义为

\[G(F)\backslash G(\A_F)\big/\big(G(\oo_{\A_F})\times K_\infty\big),\]

其中 $K_\infty$ 是 $G(F\otimes_\Q\R)$ 的某个紧子群（$F\otimes_\Q\R$ 称为 $F$ 的阿基米德完备化，它是若干 $\R$ 与 $\C$ 的拷贝的乘积）。这就是一般 $F$ 的无穷远数据。上式是一个算术局部对称空间（arithmetic locally symmetric space）。$G(F\otimes_\Q\R)$ 分解为 $G(\R)^{\times s}\times G(\C)^{\times t}$，而 $K_\infty$ 在每个因子上都要有分量。

有时人们想考虑上同调自守形式：把 $G(F)\backslash G(\A_F)/(G(\oo_{\A_F})\times K_\infty)$ 上的函数换成上同调。由于 $G_\R/K_\infty$ 可缩，这归结为群上同调

\[H^*(G(\Z);\C).\]

更一般地，若 $G(\Z)$ 作用在复向量空间 $V$ 上，可考虑 $H^*(G(\Z);V)$，即算术局部对称空间带某个局部系统扭曲后的上同调。有一个不那么显然的事实叫 Eichler–Shimura 同构（用 $\mathrm{SL}_2(\Z)\backslash\mathbb{H}$ 上的 Hodge 理论），它意味着古典模形式落在 $G=\mathrm{SL}_2$ 的上同调模形式中。与其操心带增长条件的全纯函数，不如取上同调。

函数域情形中，离散集 $\Bun_G(C)(\F_q)$ 上的函数可以等同于 $H^*(\Bun_G(C)(\F_q))$，这是自守形式理论的一个模型。一般地，「同样的」自守形式会以相当不同的函数类型出现在很不一样的地方，这是一般理论的一个怪癖。这将引向「几何自守形式」的概念，类似于「几何类域论」。我们取层：若 $k$ 是域、$C$ 是 $k$ 上曲线，如前，则考虑 $\Bun_G(C)$ 上的层，即「自守层」。这类似于阿贝尔情形：把 $\Sigma$ 送到 $\Pic(\Sigma)$ 上的层。我们希望这些层在拓扑场论中附着于曲面。

另一个可能的动机来源是数域：算术局部对称空间出现在古典数论中，所以它们的几何类比物希望正是该研究的东西。

我们遇到的基本问题是：$\Bun_G(C)$ 不是群，没有自然的群作用可用来把东西分解成特征对象。取而代之的是 Hecke 算子会作用，而令人兴奋的是它们构成交换代数，于是可以做谱分解。下次我们会更多讨论这一点。这正是场论思想大放异彩之处。
